\bookchapter{Κεφάλαιο 29: Ψηφιακή υγεία — Χτίστε $\mu$ια υγιή σχέση $\mu$ε το κινητό σας}

\chapterintro{:5 Βίντεο,?, ?,Υγεία.κεφάλαιο “”.}

\booksection{Κοιτάξτε πρώτα τον “Χρόνο οθόνης” σας — Οι αριθ$\mu$οί αφυπνίζουν}

. άνοιγ$\mu$α Λειτουργία,“Οθόνη ” (iPhone)“ψηφιακό Υγεία” (), App.

\begin{itemize}
  \item \textbf{iPhone}:Ρυθ$\mu$ίσεις > Οθόνη > Προβολή.Ξεκλείδω$\mu$α, App.
  \item \textbf{Android}:Ρυθ$\mu$ίσεις > ψηφιακό Υγείαγονείς."Υγεία ","Οθόνη".Λειτουργία, Δεδο$\mu$ένα.
\end{itemize}

\textbf{Δεδο$\mu$ένα,80\%}.4,28——.“ψηφιακό”.

\booksection{Ορίστε “όρια χρήσης” στις εφαρ$\mu$ογές — Βάλτε περιοριστή ταχύτητας στον εαυτό σας}

,Ρυθ$\mu$ίσεις περιορισ$\mu$ός.App,“”:

\begin{enumerate}
  \item \textbf{iPhone}:Ρυθ$\mu$ίσεις > Οθόνη > App > Προσθήκη.περιορισ$\mu$ός App(),Ρυθ$\mu$ίσεις 30.App, άνοιγ$\mu$α είσοδος Κωδικός.
  \item \textbf{Android}:Ρυθ$\mu$ίσεις > ψηφιακό Υγεία > Εφαρ$\mu$ογή.App Ρυθ$\mu$ίσεις, App.
  \item \textbf{Μικρή τεχνική}:Κωδικός, Κωδικός.,"".
\end{enumerate}

\booksection{Σωστή χρήση της “Μην ενοχλείτε” — Δεν κλείνεις τις ειδοποιήσεις, κλείνεις το άγχος}

Ειδοποίηση. “Ειδοποίηση ”,15-20.Μην ενοχλείτε“”.

\begin{itemize}
  \item \textbf{ενεργοποίηση}:Ρυθ$\mu$ίσεις Μην ενοχλείτε 107 αυτό$\mu$ατο Ενεργοποίηση.SMS.
  \item \textbf{επιτρέπεται}:Μην ενοχλείτε Ρυθ$\mu$ίσεις,"επιτρέπεται ""Σελιδοδείκτης Επαφή """ (3).Επείγον.
  \item \textbf{Λειτουργία εργασιακής συγκέντρωσης}:iPhone""προσαρ$\mu$οσ$\mu$ένο Σενάριο——""επιτρέπεται App Ειδοποίηση,""Ειδοποίηση Σίγαση.
\end{itemize}

\booksection{Η προστασία των $\mu$ατιών δεν εξαρτάται από εφαρ$\mu$ογές αλλά από συνήθειες — ο κανόνας 20-20-20}

Οθόνη, .(),:

\begin{tcbinfo}{Κανόνας 20-20-20}
Κάθε θέα οθόνης\textbf{20 λεπτά}, $\mu$ετακινήστε το βλέ$\mu$$\mu$α σε\textbf{20 πόδια (περίπου 6 $\mu$έτρα)}$\mu$ακριά, και κρατήστε για\textbf{20 δευτερόλεπτα}. Αυτή η απλή κίνηση ανακουφίζει αποτελεσ$\mu$ατικά την κόπωση του φωτο quiler ι$\mu$ού $\mu$υός.

\end{tcbinfo}

\begin{itemize}
  \item \textbf{ενεργοποίηση}:Ρυθ$\mu$ίσεις > ε$\mu$φάνιση > (),ενεργοποίηση Οθόνη. Συ$\mu$βουλές αυτό$\mu$ατο Ενεργοποίηση.
  \item \textbf{Προσαρ$\mu$ογή φωτεινότητας οθόνης}:Οθόνη.. ενεργοποίηση "αυτό$\mu$ατο".
  \item \textbf{Διατηρήστε απόσταση}: Το κινητό πρέπει να είναι τουλάχιστον 30 εκατοστά από τα $\mu$άτια. Πολλοί κρατάνε το κινητό πολύ κοντά, οπότε ο $\mu$υς εστιασης είναι συνεχώς σε σύσπαση και η $\mu$υωπία επιδεινώνεται γρήγορα.
\end{itemize}

\booksection{Οδηγός αυτοβοήθειας για τον τράχηλο των “σκυφτών”}

Όταν σκύβετε για να δείτε το κινητό, ο τράχηλος δέχεται πίεση 5-6 φορές $\mu$εγαλύτερη από την όρθια στάση. Η παρατετα$\mu$ένη κλίση οδηγεί σε ίσιωση της κα$\mu$πυλότητας του τραχήλου, οστικές αναπτύξεις, ζάλη και $\mu$ούδιασ$\mu$α χεριών. Ακολουθούν $\mu$ερικές απλές βελτιώσεις:

\begin{itemize}
  \item \textbf{Σηκώστε το κινητό ψηλά}: Όταν βλέπετε το κινητό, προσπαθήστε να το κρατήσετε στο ύψος των $\mu$ατιών αντί να σκύβετε για να δείτε αυτό στα γόνατα. Αν και τα χέρια θα κουραστούν, ο τράχηλος θα σας ευγνω$\mu$ονεί.
  \item \textbf{Κινήστε τον τράχηλο κάθε 30 λεπτά}: Γυρίστε αργά το κεφάλι αριστερά στο όριο, $\mu$είνετε 3 δευτερόλεπτα, επιστρέψτε. Μετά δεξιά, 3 δευτερόλεπτα. Επαναλάβετε 5 φορές. Έπειτα κάντε τη "άσκηση χαρακτήρα ρυζι " — γράψτε $\mu$ε το πηγούνι στον αέρα τον χαρακτήρα "$\mu$ι".
  \item \textbf{Σταθείτε όρθιοι στον τοίχο για 3 λεπτά}: Άκρη κεφαλιού, ω$\mu$οπλάτες, γλουτοί και πτέρνες — τέσσερα ση$\mu$εία να ακου$\mu$πούν στον τοίχο. 3 λεπτά την η$\mu$έρα διορθώνει αποτελεσ$\mu$ατικά τη στάση $\mu$ε σκυ$\mu$$\mu$ένο στήθος.
\end{itemize}

\booksection{Ύπνος και κινητό — Η συνήθεια που πρέπει να αλλάξει πρώτη στην κρεβατοκά$\mu$αρα}

Οθόνη,“”.“”, .Συ$\mu$βουλές:

\begin{itemize}
  \item \textbf{Μην κοιτάτε το κινητό 1 ώρα πριν τον ύπνο}: Αυτός είναι ο πιο αποτελεσ$\mu$ατικός αλλά και δυσκολότερος κανόνας. Αν δεν $\mu$πορείτε, τουλάχιστον ανοίξτε τη "λειτουργία νυκτός/προστασίας $\mu$ατιών" για ελάχιστο $\mu$πλε φως.
  \item \textbf{Βάλτε το κινητό εκτός της προεκ ibur ικής σας ε$\mu$βέλειας}:Φορτιστής. Ξυπνητήρι——.
  \item \textbf{ενεργοποίηση "Ύπνος"}:iPhone"Ύπνος"αυτό$\mu$ατο Ύπνος Κλείδω$\mu$α Οθόνη, Μάρκα Λειτουργία.
  \item \textbf{Ξυπνητήρι}:Ξυπνητήρι, Ξυπνητήρι,"".
\end{itemize}

\booksection{Άγχος πληροφοριών — Μάθετε την “επιλεκτική τύφλωση”}

Ο όγκος πληροφοριών που δεχό$\mu$αστε καθη$\mu$ερινά ισούται $\mu$ε αυτόν που δεχόταν οι πρόγονοί $\mu$ας σε $\mu$ια ολόκληρη ζωή. Αλλά ο ανθρώπινος εγκέφαλος δεν είναι σχεδιασ$\mu$ένος να επεξεργάζεται τόσες πληροφορίες. Η υπερφόρτωση πληροφοριών προκαλεί άγχος, διάσπαση προσοχής και κόπωση από τις αποφάσεις. Η λύση δεν είναι να “παίρνετε περισσότερες πληροφορίες ”, αλλά να “αποκλείετε ενεργητικά”:

\begin{itemize}
  \item \textbf{ακύρωση}:άνοιγ$\mu$α, UP, :"?",ακύρωση.Κανάλι.
  \item \textbf{απενεργοποίηση App}:, SMS, Η$\mu$ερολόγιο, App Ειδοποίηση.——.
  \item \textbf{Θεσπίστε $\mu$ια "η$\mu$έρα αποχής από πληροφορίες"}:(), Βίντεο..
\end{itemize}

\begin{tcbpractice}{Πρακτική: ψηφιακό Υγεία}
\begin{enumerate}
  \item :άνοιγ$\mu$α Οθόνη/ψηφιακό Υγεία, Προβολή Δεδο$\mu$ένα.
  \item :App Ρυθ$\mu$ίσεις(Συ$\mu$βουλές 30-60).
  \item :Ρυθ$\mu$ίσεις Μην ενοχλείτε ενεργοποίηση(22:00-07:00),επιτρέπεται Σελιδοδείκτης Επαφή.
  \item :1, Φορτιστής.
  \item :ακύρωση 5.
  \item :""——Βίντεο.
  \item :Προβολή Οθόνη.
\end{enumerate}

\end{tcbpractice}

\begin{tcbsummary}{Σύνοψη κεφαλαίου}
\begin{itemize}
  \item Οθόνη Δεδο$\mu$ένα,
  \item App, Μην ενοχλείτε,
  \item Ο κανόνας 20-20-20 είναι η πιο απλή και αποτελεσ$\mu$ατική συνήθεια προστασίας των $\mu$ατιών
  \item Όταν σκύβετε στο κινητό ο τράχηλος δέχεται 5-6 φορές το βάρος — σηκώστε το κινητό + κινήστε τακτικά τον τράχηλο
  \item 1, Φορτιστής
  \item Το αντίδοτο στο άγχος πληροφοριών είναι ο "ενεργητικός αποκλεισ$\mu$ός ", όχι η "περισσότερη λήψη"
\end{itemize}

\end{tcbsummary}
